% =========================================================
% Problems in Q: Sequences, Limits, and the Failure of Completeness
% =========================================================

\section{Problems in $\mathbb{Q}$}
\label{sec:problems-in-q}

\begin{tcolorbox}[
  colback=gray!6,
  colframe=gray!40,
  arc=2pt,
  left=8pt, right=8pt, top=6pt, bottom=6pt,
  title={\small\textbf{Toolkit: Problems in $\mathbb{Q}$}},
  fonttitle=\small\bfseries
]
This section assembles the formal infrastructure of rational sequences and
their limits, then uses that infrastructure to make precise exactly what goes
wrong inside $\mathbb{Q}$.  The atomic notions are:
\begin{itemize}
  \item convergence of a rational sequence to a rational limit;
  \item uniqueness of limits;
  \item subsequences and their convergence;
  \item boundedness of convergent and Cauchy sequences;
  \item closure of the Cauchy property under addition and multiplication;
  \item an equivalence relation on Cauchy sequences, which is the
        raw material for the construction of $\mathbb{R}$.
\end{itemize}
\end{tcolorbox}

% ---------------------------------------------------------
\subsection{Convergence in $\mathbb{Q}$}
\label{subsec:convergence-q}
% ---------------------------------------------------------

\begin{remark}
The rational numbers are ordered, Archimedean, and dense.  It is tempting
to believe they should already suffice for analysis.  The first serious
obstruction appears when one asks which sequences have limits, and where
those limits live.
\end{remark}

% ----- def:convergent-rational-sequence -----
\begin{tcolorbox}[colback=defbox, colframe=defborder, arc=2pt,
  left=6pt, right=6pt, top=4pt, bottom=4pt,
  title={\small\textbf{Definition (Convergence of a Rational Sequence)}},
  fonttitle=\small\bfseries]
\begin{definition}[Convergence of a rational sequence]
\label{def:convergent-rational-sequence}
Let $(x_n)$ be a sequence of rational numbers and let $a \in \mathbb{Q}$.
The sequence $(x_n)$ \emph{converges to $a$} if
\[
\forall \varepsilon \in \mathbb{Q},\; \varepsilon > 0 \;\;
\exists N \in \mathbb{N} \;\;
\forall n \geq N, \quad |x_n - a| < \varepsilon.
\]
We write $\lim_{n \to \infty} x_n = a$, or $x_n \to a$ as $n \to \infty$.
The rational number $a$ is called the \emph{limit} of $(x_n)$.
A sequence that does not converge to any element of $\mathbb{Q}$ is
called \emph{divergent in $\mathbb{Q}$}.
\end{definition}
\end{tcolorbox}

\begin{remark*}[Standard quantified statement]
\[
\forall \varepsilon > 0 \;\; \exists N \in \mathbb{N} \;\; \forall n \geq N,
\quad |x_n - a| < \varepsilon.
\]
\end{remark*}

\begin{remark*}[Definition predicate reading]
\[
\operatorname{ConvergesTo}(x_n, a)
\;\coloneqq\;
\forall \varepsilon > 0 \;\; \exists N \in \mathbb{N} \;\; \forall n \geq N,
\quad |x_n - a| < \varepsilon.
\]
\end{remark*}

\begin{remark*}[Negated quantified statement]
\[
\exists \varepsilon > 0 \;\; \forall N \in \mathbb{N} \;\; \exists n \geq N,
\quad |x_n - a| \geq \varepsilon.
\]
\end{remark*}

\begin{remark*}[Negation predicate reading]
\[
\neg\operatorname{ConvergesTo}(x_n, a)
\;\coloneqq\;
\exists \varepsilon > 0 \;\; \forall N \in \mathbb{N} \;\; \exists n \geq N,
\quad |x_n - a| \geq \varepsilon.
\]
\end{remark*}

\begin{remark*}[Failure modes]
A sequence fails to converge to $a$ in one of two ways:
\begin{enumerate}
  \item \textbf{Oscillation:} the terms return repeatedly to both sides of $a$
        with spacing bounded away from zero.  Example: $((-1)^n)$.
  \item \textbf{Unbounded drift:} the terms recede without bound and approach
        no rational number at all.  Example: $(n)$.
\end{enumerate}
\end{remark*}

\begin{remark*}[Interpretation]
Think of $\varepsilon$ as a challenge radius: the definition demands that the
sequence eventually stays inside every $\varepsilon$-ball around $a$, no
matter how small the radius.  The index $N$ is the earliest point from which
the sequence remains within that ball.  A sequence converges when it can
meet every challenge, and diverges as soon as some challenge cannot be met.
\end{remark*}

\begin{remark*}[Dependencies]
\hyperref[def:rational-sequence]{Sequence of rational numbers},
order and absolute value in $\mathbb{Q}$,
\hyperref[thm:archimedean-q]{Archimedean property of $\mathbb{Q}$}.
\end{remark*}

% ----- thm:limit-unique-q -----
\begin{theorem}[Limits of rational sequences are unique]
\label{thm:limit-unique-q}
If a sequence of rational numbers converges, then its limit is unique.
\end{theorem}

\begin{remark*}[Standard quantified statement]
\[
\forall (x_n) \subseteq \mathbb{Q},\;
\bigl(
\operatorname{ConvergesTo}(x_n, a) \land \operatorname{ConvergesTo}(x_n, b)
\bigr)
\Longrightarrow a = b.
\]
\end{remark*}

\begin{remark*}[Interpretation]
A sequence cannot approach two distinct rational numbers simultaneously.
If it did, the two alleged limits would need to be both arbitrarily close to
the same terms, forcing them to be arbitrarily close to each other — hence
equal.  This is the $2\varepsilon$ triangle-inequality squeeze.
\end{remark*}

\begin{remark*}[Dependencies]
\hyperref[def:convergent-rational-sequence]{Convergence of a rational sequence},
the triangle inequality in $\mathbb{Q}$.
\end{remark*}

\begin{remark*}[Proof]
\hyperref[prf:limit-unique-q]{Go to proof.}
\end{remark*}

% ----- def:subsequence-q -----
\begin{definition}[Subsequence of a rational sequence]
\label{def:subsequence-q}
Let $(x_n)$ be a sequence of rational numbers and let $(k_n)$ be a strictly
increasing sequence of natural numbers:
\[
k_1 < k_2 < k_3 < \cdots.
\]
The sequence $(x_{k_n})$ is called a \emph{subsequence} of $(x_n)$.
\end{definition}

\begin{remark*}[Standard quantified statement]
\[
(k_n) \subseteq \mathbb{N},\quad
\forall n \in \mathbb{N},\; k_{n+1} > k_n
\;\Longrightarrow\;
(x_{k_n}) \text{ is a subsequence of } (x_n).
\]
\end{remark*}

\begin{remark*}[Interpretation]
A subsequence is obtained by selecting an infinite subcollection of terms
in their original order — no reordering permitted, and infinitely many terms
must be retained.  The key elementary fact is that $k_n \geq n$ for all $n$,
so subsequence indices grow at least as fast as $n$ itself.
\end{remark*}

\begin{remark*}[Dependencies]
\hyperref[def:rational-sequence]{Sequence of rational numbers}.
\end{remark*}

% ----- thm:subsequence-convergence-q -----
\begin{theorem}[Subsequences of convergent sequences converge to the same limit]
\label{thm:subsequence-convergence-q}
Let $(x_n)$ be a sequence of rational numbers with $x_n \to a$.
Then every subsequence $(x_{k_n})$ of $(x_n)$ also converges to $a$:
\[
\lim_{n \to \infty} x_{k_n} = a.
\]
\end{theorem}

\begin{remark*}[Standard quantified statement]
\[
\operatorname{ConvergesTo}(x_n, a)
\;\Longrightarrow\;
\forall (k_n) \text{ strictly increasing},\;
\operatorname{ConvergesTo}(x_{k_n}, a).
\]
\end{remark*}

\begin{remark*}[Contrapositive quantified statement]
\[
\exists (k_n) \text{ strictly increasing},\;
\neg\operatorname{ConvergesTo}(x_{k_n}, a)
\;\Longrightarrow\;
\neg\operatorname{ConvergesTo}(x_n, a).
\]
\end{remark*}

\begin{remark*}[Interpretation]
The contrapositive is the working tool: to prove that $(x_n)$ diverges it
suffices to find two subsequences converging to distinct limits, or one
subsequence that diverges.  For example, $((-1)^n)$ has subsequences
$((-1)^{2n}) = (1)$ and $((-1)^{2n+1}) = (-1)$, with limits $1 \neq -1$,
so $((-1)^n)$ diverges.
\end{remark*}

\begin{remark*}[Dependencies]
\hyperref[def:convergent-rational-sequence]{Convergence of a rational sequence},
\hyperref[def:subsequence-q]{Subsequence}.
\end{remark*}

\begin{remark*}[Proof]
\hyperref[prf:subsequence-convergence-q]{Go to proof.}
\end{remark*}

% ---------------------------------------------------------
\subsection{Boundedness}
\label{subsec:boundedness-sequences-q}
% ---------------------------------------------------------

\begin{remark}
Boundedness is a coarser property than convergence, but it is entailed by
convergence and also by the Cauchy property.  These implications are one-way:
the sequence $((-1)^n)$ is bounded but divergent, so boundedness does not
imply convergence.
\end{remark}

% ----- def:bounded-sequence-q -----
\begin{definition}[Bounded sequence of rational numbers]
\label{def:bounded-sequence-q}
A sequence $(x_n)$ of rational numbers is:
\begin{itemize}
  \item \emph{bounded above} if $\exists \Lambda \in \mathbb{Q}\; \forall n \in \mathbb{N},\;
        x_n \leq \Lambda$;
  \item \emph{bounded below} if $\exists \lambda \in \mathbb{Q}\; \forall n \in \mathbb{N},\;
        x_n \geq \lambda$;
  \item \emph{bounded} if it is bounded above and bounded below, equivalently if
        $\exists \Gamma \in \mathbb{Q}\; \forall n \in \mathbb{N},\; |x_n| \leq \Gamma$.
\end{itemize}
\end{definition}

\begin{remark*}[Standard quantified statement]
\[
\operatorname{Bounded}(x_n)
\;\coloneqq\;
\exists \Gamma \in \mathbb{Q},\; \Gamma > 0,\;\;
\forall n \in \mathbb{N},\;\; |x_n| \leq \Gamma.
\]
\end{remark*}

\begin{remark*}[Negated quantified statement]
\[
\neg\operatorname{Bounded}(x_n)
\;\coloneqq\;
\forall \Gamma > 0 \;\; \exists n \in \mathbb{N},\;\; |x_n| > \Gamma.
\]
\end{remark*}

\begin{remark*}[Interpretation]
Boundedness is a global envelope condition: one single rational bound must
contain every term.  The absolute-value formulation is equivalent to
specifying both an upper bound $\Gamma$ and a lower bound $-\Gamma$.
\end{remark*}

\begin{remark*}[Dependencies]
\hyperref[def:rational-sequence]{Sequence of rational numbers},
absolute value in $\mathbb{Q}$.
\end{remark*}

% ----- thm:convergent-bounded-q -----
\begin{theorem}[Convergent sequences are bounded]
\label{thm:convergent-bounded-q}
Every convergent sequence of rational numbers is bounded.
\end{theorem}

\begin{remark*}[Standard quantified statement]
\[
\operatorname{ConvergesTo}(x_n, a)
\;\Longrightarrow\;
\operatorname{Bounded}(x_n).
\]
\end{remark*}

\begin{remark*}[Interpretation]
Convergence forces the tail of the sequence into a bounded neighborhood of
the limit; the finite initial segment up to index $N$ is bounded by the
maximum of finitely many terms.  Together these cover all indices.
The converse fails: $((-1)^n)$ is bounded but does not converge.
\end{remark*}

\begin{remark*}[Dependencies]
\hyperref[def:convergent-rational-sequence]{Convergence},
\hyperref[def:bounded-sequence-q]{Bounded sequence},
the triangle inequality and Archimedean property of $\mathbb{Q}$.
\end{remark*}

\begin{remark*}[Proof]
\hyperref[prf:convergent-bounded-q]{Go to proof.}
\end{remark*}

% ----- thm:cauchy-bounded-q -----
\begin{theorem}[Cauchy sequences are bounded]
\label{thm:cauchy-bounded-q}
Every Cauchy sequence of rational numbers is bounded.
\end{theorem}

\begin{remark*}[Standard quantified statement]
\[
\operatorname{CauchySequence}(x_n)
\;\Longrightarrow\;
\operatorname{Bounded}(x_n).
\]
\end{remark*}

\begin{remark*}[Interpretation]
The Cauchy condition forces the tail of the sequence — all terms beyond
index $N$ — to lie within distance $1$ of the single term $x_{N+1}$.
That tail is therefore bounded by $|x_{N+1}| + 1$.  The finite initial
segment contributes finitely many terms whose absolute values can be
bounded by their maximum.
\end{remark*}

\begin{remark*}[Dependencies]
\hyperref[def:cauchy-rational]{Cauchy sequence in $\mathbb{Q}$},
\hyperref[def:bounded-sequence-q]{Bounded sequence}.
\end{remark*}

\begin{remark*}[Proof]
\hyperref[prf:cauchy-bounded-q]{Go to proof.}
\end{remark*}

% ---------------------------------------------------------
\subsection{Algebraic Closure of Cauchy Sequences}
\label{subsec:cauchy-algebra-q}
% ---------------------------------------------------------

\begin{remark}
The Cauchy property is not merely stable; it is closed under the algebraic
operations of $\mathbb{Q}$.  This is essential for the Cauchy-sequence
construction of $\mathbb{R}$: one needs to add and multiply equivalence
classes of Cauchy sequences, so the Cauchy property must survive these
operations.
\end{remark}

% ----- thm:cauchy-closed-operations-q -----
\begin{theorem}[Cauchy sequences are closed under addition and multiplication]
\label{thm:cauchy-closed-operations-q}
If $(x_n)$ and $(y_n)$ are Cauchy sequences of rational numbers, then so
are $(x_n + y_n)$ and $(x_n y_n)$.
\end{theorem}

\begin{remark*}[Standard quantified statement]
\[
\operatorname{CauchySequence}(x_n) \land \operatorname{CauchySequence}(y_n)
\;\Longrightarrow\;
\operatorname{CauchySequence}(x_n + y_n)
\;\land\;
\operatorname{CauchySequence}(x_n y_n).
\]
\end{remark*}

\begin{remark*}[Interpretation]
For addition the argument is a direct $\varepsilon/2$ split.
For multiplication the key is that Cauchy sequences are bounded
(\hyperref[thm:cauchy-bounded-q]{Theorem~\ref*{thm:cauchy-bounded-q}}),
so one can bound the cross-terms $|y_n||x_n - x_m|$ and $|x_m||y_n - y_m|$
using a common bound $\Lambda$ for both sequences, then split $\varepsilon$
as $\varepsilon/(2\Lambda)$ for each piece.
\end{remark*}

\begin{remark*}[Dependencies]
\hyperref[def:cauchy-rational]{Cauchy sequence in $\mathbb{Q}$},
\hyperref[thm:cauchy-bounded-q]{Cauchy sequences are bounded},
the triangle inequality in $\mathbb{Q}$.
\end{remark*}

\begin{remark*}[Proof]
\hyperref[prf:cauchy-closed-operations-q]{Go to proof.}
\end{remark*}

% ---------------------------------------------------------
\subsection{Equivalence of Cauchy Sequences}
\label{subsec:cauchy-equivalence-q}
% ---------------------------------------------------------

\begin{remark}
The Cauchy-sequence construction of $\mathbb{R}$ does not identify $\mathbb{R}$
with individual Cauchy sequences but with equivalence classes of them.
Two Cauchy sequences are equivalent if they have the same ``eventual
behavior'' — that is, if their difference converges to zero.
This section introduces that equivalence relation and verifies that it
is indeed an equivalence relation, which is a prerequisite for forming
quotient sets.
\end{remark}

% ----- def:cauchy-equivalence -----
\begin{tcolorbox}[colback=defbox, colframe=defborder, arc=2pt,
  left=6pt, right=6pt, top=4pt, bottom=4pt,
  title={\small\textbf{Definition (Equivalence of Cauchy Sequences in $\mathbb{Q}$)}},
  fonttitle=\small\bfseries]
\begin{definition}[Equivalence of Cauchy sequences in $\mathbb{Q}$]
\label{def:cauchy-equivalence}
Let $\mathbb{Q}^c$ denote the set of all Cauchy sequences of rational numbers.
Two sequences $(x_n), (y_n) \in \mathbb{Q}^c$ are \emph{equivalent}, written
$(x_n) \sim (y_n)$, if
\[
\forall \varepsilon > 0 \;\; \exists N \in \mathbb{N} \;\; \forall n \geq N,
\quad |x_n - y_n| < \varepsilon.
\]
Equivalently, $(x_n) \sim (y_n)$ if and only if $(x_n - y_n) \sim (0)$,
where $(0)$ denotes the constant zero sequence.
\end{definition}
\end{tcolorbox}

\begin{remark*}[Standard quantified statement]
\[
(x_n) \sim (y_n)
\;\coloneqq\;
\forall \varepsilon > 0 \;\; \exists N \in \mathbb{N} \;\; \forall n \geq N,
\quad |x_n - y_n| < \varepsilon.
\]
\end{remark*}

\begin{remark*}[Definition predicate reading]
\[
\operatorname{CauchyEquiv}(x_n, y_n)
\;\coloneqq\;
\operatorname{ConvergesTo}(x_n - y_n,\, 0).
\]
Two Cauchy sequences are equivalent exactly when their pointwise difference
is a null sequence.
\end{remark*}

\begin{remark*}[Negated quantified statement]
\[
\neg\bigl((x_n) \sim (y_n)\bigr)
\;\coloneqq\;
\exists \varepsilon > 0 \;\; \forall N \in \mathbb{N} \;\; \exists n \geq N,
\quad |x_n - y_n| \geq \varepsilon.
\]
\end{remark*}

\begin{remark*}[Interpretation]
Two Cauchy sequences are equivalent when they are eventually indistinguishable
at every scale.  Geometrically they converge to the same ``ideal point,''
even if that point has no rational representative.  The equivalence classes
under $\sim$ are exactly the real numbers in the Cauchy-sequence construction
of $\mathbb{R}$.
\end{remark*}

\begin{remark*}[Dependencies]
\hyperref[def:cauchy-rational]{Cauchy sequence in $\mathbb{Q}$},
\hyperref[def:convergent-rational-sequence]{Convergence of a rational sequence}.
\end{remark*}

% ----- thm:cauchy-equivalence-is-equivalence -----
\begin{tcolorbox}[colback=thmbox, colframe=thmborder, arc=2pt,
  left=6pt, right=6pt, top=4pt, bottom=4pt,
  title={\small\textbf{Theorem (The Cauchy Equivalence Relation on $\mathbb{Q}^c$)}},
  fonttitle=\small\bfseries]
\begin{theorem}[The relation $\sim$ is an equivalence relation on $\mathbb{Q}^c$]
\label{thm:cauchy-equivalence-is-equivalence}
The relation $\sim$ defined on $\mathbb{Q}^c$ by
$(x_n) \sim (y_n) \Longleftrightarrow \operatorname{ConvergesTo}(x_n - y_n, 0)$
is an equivalence relation.
\end{theorem}
\end{tcolorbox}

\begin{remark*}[Standard quantified statement]
The relation $\sim$ on $\mathbb{Q}^c$ satisfies:
\begin{align*}
&\forall (x_n) \in \mathbb{Q}^c,\quad (x_n) \sim (x_n); \\
&\forall (x_n),(y_n) \in \mathbb{Q}^c,\quad
  (x_n) \sim (y_n) \Longrightarrow (y_n) \sim (x_n); \\
&\forall (x_n),(y_n),(z_n) \in \mathbb{Q}^c,\quad
  (x_n) \sim (y_n) \land (y_n) \sim (z_n)
  \Longrightarrow (x_n) \sim (z_n).
\end{align*}
\end{remark*}

\begin{remark*}[Interpretation]
Reflexivity is immediate: $|x_n - x_n| = 0 < \varepsilon$ always.
Symmetry follows because $|x_n - y_n| = |y_n - x_n|$.
Transitivity is the triangle inequality applied to $|x_n - z_n| \leq
|x_n - y_n| + |y_n - z_n|$ with an $\varepsilon/2$ split — identical in
structure to the proof that convergent sequences have unique limits.
\end{remark*}

\begin{remark*}[Dependencies]
\hyperref[def:cauchy-equivalence]{Equivalence of Cauchy sequences},
\hyperref[thm:cauchy-closed-operations-q]{Cauchy sequences closed under addition},
triangle inequality in $\mathbb{Q}$.
\end{remark*}

\begin{remark*}[Proof]
\hyperref[prf:cauchy-equivalence-is-equivalence]{Go to proof.}
\end{remark*}

% ----- thm:cauchy-equiv-preserved-under-operations -----
\begin{theorem}[Cauchy equivalence is preserved under addition and multiplication]
\label{thm:cauchy-equiv-preserved-operations}
Let $(x_n),(y_n),(u_n),(v_n) \in \mathbb{Q}^c$.
Suppose $(x_n) \sim (u_n)$ and $(y_n) \sim (v_n)$.  Then:
\[
(x_n + y_n) \sim (u_n + v_n)
\qquad\text{and}\qquad
(x_n y_n) \sim (u_n v_n).
\]
\end{theorem}

\begin{remark*}[Standard quantified statement]
\[
(x_n) \sim (u_n) \land (y_n) \sim (v_n)
\;\Longrightarrow\;
(x_n + y_n) \sim (u_n + v_n)
\;\land\;
(x_n y_n) \sim (u_n v_n).
\]
\end{remark*}

\begin{remark*}[Interpretation]
This is the well-definedness result for arithmetic on equivalence classes.
Without it, adding or multiplying representatives from two classes could
land in different classes depending on which representatives were chosen,
making the operations on $\mathbb{R} = \mathbb{Q}^c / \sim$ ill-defined.
The proof uses boundedness of Cauchy sequences (\hyperref[thm:cauchy-bounded-q]{Theorem~\ref*{thm:cauchy-bounded-q}})
in the same way that the analogous result for multiplication on $\mathbb{Q}$
used boundedness of the fraction representatives.
\end{remark*}

\begin{remark*}[Dependencies]
\hyperref[def:cauchy-equivalence]{Cauchy equivalence},
\hyperref[thm:cauchy-closed-operations-q]{Cauchy sequences closed under operations},
\hyperref[thm:cauchy-bounded-q]{Cauchy sequences are bounded}.
\end{remark*}

\begin{remark*}[Proof]
\hyperref[prf:cauchy-equiv-preserved-operations]{Go to proof.}
\end{remark*}

% ---------------------------------------------------------
\subsection{The Precise Statement of Incompleteness}
\label{subsec:incompleteness-statement-q}
% ---------------------------------------------------------

\begin{remark}
With convergence and the Cauchy property now formally in place, the
incompleteness of $\mathbb{Q}$ can be stated with full precision.
\end{remark}

\begin{remark}[What the definitions make precise]
\hyperref[def:convergent-rational-sequence]{Convergence} requires
a limit \emph{in $\mathbb{Q}$}.
\hyperref[def:cauchy-rational]{The Cauchy property} requires only that
the terms become close to \emph{each other}.
These are different conditions.
\hyperref[thm:convergent-implies-cauchy-q]{Theorem~\ref*{thm:convergent-implies-cauchy-q}}
shows convergence implies Cauchy.
\hyperref[thm:cauchy-not-convergent-q]{Theorem~\ref*{thm:cauchy-not-convergent-q}}
shows the converse fails.
That failure is incompleteness.
\end{remark}

\begin{remark}[Diagram of implications in $\mathbb{Q}$]
\[
\text{convergent in } \mathbb{Q}
\;\Longrightarrow\;
\text{Cauchy in } \mathbb{Q}
\;\Longrightarrow\;
\text{bounded in } \mathbb{Q},
\]
\[
\text{bounded}
\;\not\Longrightarrow\;
\text{Cauchy},
\qquad
\text{Cauchy}
\;\not\Longrightarrow\;
\text{convergent in } \mathbb{Q}.
\]
Each failure is witnessed by an explicit counterexample.
The last implication holds in $\mathbb{R}$ but not in $\mathbb{Q}$:
that is exactly the completeness axiom.
\end{remark}
